\documentclass[11pt,letterpaper]{article}
\usepackage[T1]{fontenc}
\usepackage{lmodern}
\usepackage[margin=1in]{geometry}
\usepackage{amsmath,amssymb,amsthm,mathtools}
\usepackage{booktabs,longtable,array}
\usepackage{microtype}
\usepackage{xcolor}
\usepackage{hyperref}
\usepackage{fancyhdr}
\usepackage{listings}
\usepackage{enumitem}
\usepackage{graphicx}

\definecolor{inkblue}{RGB}{27,57,88}
\definecolor{quietgray}{RGB}{85,85,85}
\hypersetup{
  colorlinks=true,
  linkcolor=inkblue,
  citecolor=inkblue,
  urlcolor=inkblue,
  pdftitle={Winding Sectors and the Four-Displacement Restricted-Permutation Problem},
  pdfauthor={Research report for the ProveIt project}
}
\pagestyle{fancy}
\fancyhf{}
\fancyhead[L]{\small\textsc{Four-displacement restricted permutations}}
\fancyhead[R]{\small October 2026}
\fancyfoot[C]{\thepage}
\setlength{\headheight}{14pt}
\setlength{\parskip}{4pt}
\setlength{\parindent}{1.2em}
\setcounter{tocdepth}{2}
\numberwithin{equation}{section}
\allowdisplaybreaks

\newtheorem{theorem}{Theorem}[section]
\newtheorem{lemma}[theorem]{Lemma}
\newtheorem{proposition}[theorem]{Proposition}
\newtheorem{corollary}[theorem]{Corollary}
\theoremstyle{definition}
\newtheorem{definition}[theorem]{Definition}
\newtheorem{remark}[theorem]{Remark}
\newtheorem{question}[theorem]{Research question}

\newcommand{\Z}{\mathbb Z}
\newcommand{\N}{\mathbb N}
\newcommand{\C}{\mathbb C}
\newcommand{\Zn}{\Z/n\Z}
\newcommand{\per}{\operatorname{per}}
\newcommand{\GF}{\mathcal G}
\newcommand{\W}{\mathcal W}
\newcommand{\calC}{\mathcal C}
\newcommand{\bigO}{\mathcal O}
\newcommand{\e}{\mathrm e}
\newcommand{\ii}{\mathrm i}
\newcommand{\oeis}[1]{\texttt{#1}}
\newcommand{\abs}[1]{\lvert#1\rvert}
\newcommand{\floor}[1]{\left\lfloor#1\right\rfloor}
\newcommand{\ceil}[1]{\left\lceil#1\right\rceil}
\DeclareMathOperator*{\argmin}{arg\,min}

\lstset{
  basicstyle=\small\ttfamily,
  breaklines=true,
  columns=fullflexible,
  frame=single,
  rulecolor=\color{quietgray},
  xleftmargin=0.5em,
  xrightmargin=0.5em,
  showstringspaces=false
}

\title{\textbf{Winding Sectors and the Four-Displacement\\Restricted-Permutation Problem}\\[0.7em]
\large A direct Tribonacci--Lucas enumeration, an inverse transseries,\\
and corrections to OEIS A000382, A000496, and A008305}
\author{Research report for the ProveIt project}
\date{October 1, 2026}

\begin{document}
\maketitle

\begin{abstract}
For $n\ge4$, let $\Phi_n$ be the number of permutations $\sigma$ of
$\Z/n\Z$ for which every clockwise displacement
$\sigma(i)-i\pmod n$ belongs to $\{0,1,2,3\}$.  Equivalently,
$\Phi_n=\per(I+P+P^2+P^3)$ for the cyclic shift matrix $P$.  This is the
fourth column of OEIS A008305 and the tail of A004306.  Current OEIS entries
also identify it with A000496 and with four times A000382, but those sequences
first disagree at $n=8$.

We give a short direct proof of the exact formula
\[
  \Phi_n=2(L_n+1),
  \qquad L_0=3,\ L_1=1,\ L_2=3,
  \qquad L_n=L_{n-1}+L_{n-2}+L_{n-3}.
\]
The proof partitions every allowed permutation by its integer winding number.
A constant-cut lemma shows that winding-one permutations are precisely cyclic
tilings by monomers, dimers, and trimers; reflection-complementation identifies
the winding-two sector with the winding-one sector, while the two outer
sectors are singletons.  The method extends to the outer winding sectors for
all displacement widths.

Consequences include the corrected recurrence and rational generating
function, an exact Binet formula, exponentially sharp asymptotics, a
nearest-integer formula, and an all-orders logarithmically oscillatory inverse
transseries.  We also prove the recurrence currently marked ``conjectured''
in A000382, but for a different reason: it is an identity for the sequence
produced by Mendelsohn's displayed 1961 formula, not for the restricted-
permutation count.  The exact correction is
\[
 \sum_{n\ge4}\left(\frac{\Phi_n}{4}-A000382(n)\right)z^{n-4}
 =\frac{z^4}{1-2z+z^4}.
\]
An independent subset-dynamic-programming verifier is included.  This is a
research draft; the historical and novelty claims should be independently
checked before an OEIS or journal submission.
\end{abstract}

\clearpage
{\small\setlength{\parskip}{0pt}\tableofcontents}
\clearpage

\section{The problem, the discrepancy, and the result}

\subsection{The cyclic restricted-permutation count}

Write the elements of $\Zn$ as $0,1,\dots,n-1$.  For a permutation
$\sigma\colon\Zn\to\Zn$, let $d_i$ be the unique integer in
$\{0,1,\dots,n-1\}$ satisfying
\[
  \sigma(i)=i+d_i\pmod n.
\]
For $n\ge4$ define
\begin{equation}\label{eq:Phi-definition}
  \Phi_n
  :=\#\bigl\{\sigma\in S(\Zn):d_i\in\{0,1,2,3\}
       \text{ for every }i\bigr\}.
\end{equation}
If $P$ is the permutation matrix for the clockwise cyclic shift, then
\begin{equation}\label{eq:permanent}
  \Phi_n=\per(I+P+P^2+P^3).
\end{equation}
This is Mendelsohn's $\phi(n,4)$ in the notation of his 1961 paper
\cite{mendelsohn1961}; it is also $T(n,4)$ in OEIS A008305
\cite{oeisA008305}.

The first values obtained from the definition are
\begin{equation}\label{eq:true-first-values}
  24,44,80,144,264,484,888,1632,3000,5516,10144,\dots
  \qquad(n=4,5,6,\dots).
\end{equation}
They are recorded in OEIS A004306, whose current entry gives the recurrence
$a(n)=2a(n-1)-a(n-4)$ and a Tribonacci formula
\cite{oeisA004306}.  The purpose of this report is not to rediscover that
recorded recurrence by sequence fitting.  It is to give a transparent
bijection explaining it, to separate three historically entangled sequences,
and to derive several consequences that do not appear in the checked entries.

\subsection{Three sequences that agree only through n=7}

As of October 1, 2026, OEIS A000382 begins
\[
 6,11,20,36,65,119,218,400,735,\dots
\]
with offset $4$, and says that it is the fourth column of A008305 divided by
$4$; it also marks
\[
 a(n)=a(n-1)+a(n-2)+a(n-3)-2
\]
as conjectural \cite{oeisA000382}.  A000496 begins
\[
 1,1,2,6,24,44,80,144,260,476,872,\dots
\]
and is likewise named ``Restricted permutations'' \cite{oeisA000496}.
However, A008305 itself gives $264,484,888$ in the relevant column at
$n=8,9,10$, and A004306 gives the same values
\cite{oeisA008305,oeisA004306}.

The divergence is summarized below.  ``Equation (4)'' means the sequence
obtained by evaluating Mendelsohn's displayed formula, while ``printed table''
means the values printed on page 32 of the same paper.

\begin{center}
\begin{tabular}{c@{\quad}rrrrrrr}
\toprule
$n$ & 4 & 5 & 6 & 7 & 8 & 9 & 10\\
\midrule
actual $\Phi_n$ & 24 & 44 & 80 & 144 & 264 & 484 & 888\\
Mendelsohn equation (4) & 24 & 44 & 80 & 144 & 260 & 476 & 872\\
Mendelsohn printed table & 24 & 44 & 80 & 144 & 256 & 472 & 872\\
$4\,A000382(n)=A000496(n)$ & 24 & 44 & 80 & 144 & 260 & 476 & 872\\
\bottomrule
\end{tabular}
\end{center}

Thus the source formula, the source table, and the permanent are three distinct
objects.  The first discrepancy between the current A000382 comment and the
actual fourth column occurs at
\[
  n=8:\qquad A000382(8)=65,
  \qquad \frac{T(8,4)}4=\frac{264}{4}=66.
\]
The issue is not a matter of offset convention.

\subsection{Main theorem}

Let $L_n$ denote the Tribonacci--Lucas sequence
\begin{equation}\label{eq:Lucas-def}
  L_0=3,\qquad L_1=1,\qquad L_2=3,
  \qquad L_n=L_{n-1}+L_{n-2}+L_{n-3}.
\end{equation}
This is OEIS A001644 \cite{oeisA001644}.

\begin{theorem}[Exact four-displacement enumeration]\label{thm:main}
For every $n\ge4$,
\begin{equation}\label{eq:main-formula}
  \boxed{\Phi_n=2(L_n+1).}
\end{equation}
More precisely, if the allowed permutations are sorted by winding number
$q\in\{0,1,2,3\}$, then the four sector sizes are
\begin{equation}\label{eq:sector-vector}
  (1,L_n,L_n,1).
\end{equation}
\end{theorem}

The proof occupies Sections~\ref{sec:winding}--\ref{sec:specialization}.
It uses only a conservation law and a cyclic-tiling bijection.

\subsection{What is new here, and what is not claimed}

The recurrence and generating function for the correct count are already
present in A004306, and related transfer-matrix formulae occur in the classical
literature on permanents of cyclic $(0,1)$ matrices
\cite{metropolis1969,beker1987}.  The contributions of this report are:

\begin{enumerate}[label=(\roman*),leftmargin=2.4em]
\item a direct winding-sector proof of the exact formula;
\item a general theorem identifying the outer winding sectors for every fixed
      displacement width;
\item a source-level diagnosis of the A000382/A000496 discrepancy;
\item a proof of the A000382 recurrence for the formula-defined, rather than
      permutation-defined, sequence;
\item the exact correction kernel and its asymptotics;
\item exact and all-orders asymptotic formulae, including an inverse
      log-periodic transseries.
\end{enumerate}

No claim is made that no equivalent bijection has ever appeared in the wider
literature.  The literature search was targeted rather than exhaustive.  The
OEIS statements quoted here are snapshots and may be corrected after the date
of this report.

\section{Winding sectors and a constant-cut law}\label{sec:winding}

It is useful to treat a slightly more general problem.  Fix integers
$n\ge r\ge2$, and define
\begin{equation}\label{eq:Phi-nr}
  \Phi_{n,r}
  :=\#\bigl\{\sigma\in S(\Zn):
       \sigma(i)-i\pmod n\in\{0,1,\dots,r-1\}\bigr\}.
\end{equation}
For every allowed permutation choose the representatives
$d_i\in\{0,1,\dots,r-1\}$.

\begin{lemma}[Integral winding number]\label{lem:winding-integral}
For every allowed permutation,
\begin{equation}\label{eq:q-def}
  q(\sigma):=\frac1n\sum_{i=0}^{n-1}d_i
\end{equation}
is an integer in $\{0,1,\dots,r-1\}$.
\end{lemma}

\begin{proof}
Because $\sigma$ is a permutation,
\[
 \sum_i\sigma(i)\equiv\sum_i i\pmod n.
\]
Hence $\sum_i d_i\equiv0\pmod n$.  The bounds
$0\le d_i\le r-1$ give $0\le q\le r-1$.
\end{proof}

Draw the vertices $0,1,\dots,n-1$ clockwise on a circle.  Represent the arrow
$i\mapsto i+d_i$ by the clockwise arc that traverses exactly $d_i$ elementary
cuts.  Let cut $k$ be the gap from $k$ to $k+1$, and let $c_k$ be the number
of nontrivial arcs crossing that cut.

\begin{lemma}[Constant-cut law]\label{lem:constant-cut}
Every cut is crossed by exactly $q(\sigma)$ arcs:
\begin{equation}\label{eq:constant-cut}
  c_0=c_1=\cdots=c_{n-1}=q(\sigma).
\end{equation}
\end{lemma}

\begin{proof}
Moving from cut $k-1$ to cut $k$, the crossing count gains one if a nonloop
arc starts at $k$ and loses one if a nonloop arc ends at $k$.  There is one
outgoing and one incoming arrow at every vertex.  Moreover, the outgoing arrow
at $k$ is a loop if and only if the incoming arrow at $k$ is that same loop:
if $\sigma(k)=k$, no other arrow can end at $k$, and conversely.  Therefore the
number of nonloop starts at $k$ equals the number of nonloop ends at $k$, so
$c_k-c_{k-1}=0$.

The sum of the $n$ cut counts is the total number of elementary cuts traversed
by all arcs, namely $\sum_i d_i=qn$.  Since the cut counts are constant, each
is $q$.
\end{proof}

This elementary observation is the structural core of the argument.  It turns
a modular displacement condition into a covering problem on a circle.

\subsection{Reflection-complementation}

Define an operation on allowed permutations by
\begin{equation}\label{eq:complement-map}
  (\calC_r\sigma)(x):=r-1-\sigma(-x)\pmod n.
\end{equation}
It is the composition of reflection in the domain, the permutation $\sigma$,
and an affine reflection in the codomain.

\begin{lemma}[Sector symmetry]\label{lem:sector-symmetry}
The map $\calC_r$ is an involution, and it sends winding sector $q$ to winding
sector $r-1-q$.  In particular, the winding-sector polynomial
\begin{equation}\label{eq:sector-polynomial}
  W_{n,r}(u):=\sum_{\sigma}u^{q(\sigma)}
\end{equation}
is palindromic of degree $r-1$.
\end{lemma}

\begin{proof}
If $d_{-x}$ is the displacement of $\sigma$ at $-x$, then
\[
 (\calC_r\sigma)(x)-x
 \equiv r-1-\sigma(-x)-x
 \equiv r-1-d_{-x}\pmod n.
\]
The displayed integer already lies in $\{0,\dots,r-1\}$.  Summing over $x$
changes $qn$ into $(r-1)n-qn$.  A second application of $\calC_r$ returns
$\sigma$.
\end{proof}

The extreme sectors are immediate.

\begin{lemma}[Outer singleton sectors]\label{lem:outer-singletons}
The winding-$0$ sector consists only of the identity, and the winding-$(r-1)$
sector consists only of the shift $i\mapsto i+r-1$.
\end{lemma}

\begin{proof}
An average of numbers in $[0,r-1]$ equals $0$ only when every number is $0$,
and it equals $r-1$ only when every number is $r-1$.
\end{proof}

\section{Winding one is a cyclic-tiling problem}\label{sec:tilings}

Let $k=r-1$.  A cyclic tiling of a labeled $n$-cycle by tiles of lengths
$1,2,\dots,k$ is a partition of its elementary cuts into clockwise intervals
of those lengths.  A tile has a distinguished initial vertex, so the labels
and orientation are retained.

\begin{theorem}[Outer-sector tiling theorem]\label{thm:outer-sector}
For $n\ge r$, the winding-$1$ permutations counted by $\Phi_{n,r}$ are in
bijection with cyclic tilings of a labeled $n$-cycle by tiles of lengths
$1,2,\dots,r-1$.

Under the bijection, a tile of length $j$ becomes the displacement block
\begin{equation}\label{eq:tile-block-general}
  j\,\underbrace{00\cdots0}_{j-1\text{ zeros}}.
\end{equation}
The winding-$(r-2)$ sector has the same size by
Lemma~\ref{lem:sector-symmetry}.
\end{theorem}

\begin{proof}
By Lemma~\ref{lem:constant-cut}, every cut is covered by exactly one arc in a
winding-$1$ permutation.  Consequently, the cut intervals traversed by its
nonloop arcs are pairwise disjoint and cover the cycle.  If an arc starts at
$i$ and has length $j$, then no nonloop arc can start at an internal vertex
$i+1,\dots,i+j-1$, because its first cut would already be covered.  Those
internal vertices are therefore loops, giving the block in
\eqref{eq:tile-block-general}.  The nonloop arcs thus define a cyclic tiling.

Conversely, replace each tile of length $j$ by an arc from its initial vertex
to its terminal vertex, and put loops on the $j-1$ internal vertices.  Every
internal vertex is its own target, while every tile initial vertex is the
target of the preceding tile's arc.  Thus every vertex has exactly one
preimage and the construction is a permutation.  Each cut is crossed once, so
its winding number is $1$.  The two constructions are inverse.
\end{proof}

For $r=4$, the three tile blocks are particularly simple:
\begin{equation}\label{eq:blocks-123}
  1,\qquad 20,\qquad 300.
\end{equation}
For example, the cyclic composition $8=3+2+1+2$ gives the displacement word
\[
  300\,20\,1\,20=30020120.
\]

\subsection{Counting the cyclic tilings}

Let $F_m^{(k)}$ be the number of linear tilings of length $m$ with tile lengths
$1,\dots,k$.  Set $F_0^{(k)}=1$ and $F_m^{(k)}=0$ for $m<0$.  Then
\begin{equation}\label{eq:kbonacci-compositions}
  F_m^{(k)}=\sum_{j=1}^k F_{m-j}^{(k)},
  \qquad
  \sum_{m\ge0}F_m^{(k)}z^m=\frac1{1-z-z^2-\cdots-z^k}.
\end{equation}

\begin{proposition}[Cyclic-tiling count]\label{prop:cyclic-count}
Let $C_{n,k}$ be the number of cyclic tilings of a labeled $n$-cycle with tile
lengths $1,\dots,k$.  Then, for $n\ge1$,
\begin{equation}\label{eq:C-general}
  C_{n,k}=\sum_{j=1}^k jF_{n-j}^{(k)}.
\end{equation}
It satisfies the same order-$k$ recurrence as $F_m^{(k)}$ once the initial
range has passed.  Equivalently, if $\lambda_1,\dots,\lambda_k$ are the roots
of
\begin{equation}\label{eq:k-characteristic}
  x^k-x^{k-1}-\cdots-x-1=0,
\end{equation}
then
\begin{equation}\label{eq:power-sum-general}
  C_{n,k}=\lambda_1^n+\cdots+\lambda_k^n.
\end{equation}
\end{proposition}

\begin{proof}
Focus on the unique tile containing vertex $0$.  If that tile has length $j$,
then vertex $0$ can occupy any of its $j$ labeled positions.  Removing it
leaves a line of length $n-j$, which can be tiled in $F_{n-j}^{(k)}$ ways.
This proves \eqref{eq:C-general}.

The power sums in \eqref{eq:power-sum-general} obey the recurrence associated
with \eqref{eq:k-characteristic}.  Newton's identities give the same initial
values as \eqref{eq:C-general}; alternatively, both sides are the trace of the
$n$th power of the standard composition automaton.  Hence they agree.
\end{proof}

The dominant root $\lambda_k$ of \eqref{eq:k-characteristic} is real, simple,
and lies in $(1,2)$.  Multiplication by $x-1$ shows that it is the nontrivial
root in $(1,2)$ of
\begin{equation}\label{eq:Mendelsohn-root-general}
  x^{k+1}-2x^k+1=0.
\end{equation}
Thus the outer winding sector already explains why generalized Fibonacci roots
appear in confined-displacement problems.

\section{Complete solution for four displacements}\label{sec:specialization}

Set $r=4$ and $k=3$.  The possible winding numbers are $0,1,2,3$.  There are no
unaccounted-for inner sectors: winding $0$ and $3$ are singletons, winding $1$
is the cyclic-tiling sector, and winding $2$ is its reflected complement.

Write simply $F_m=F_m^{(3)}$.  Thus
\begin{equation}\label{eq:F-initial}
  F_0=1,\quad F_1=1,\quad F_2=2,\quad
  F_m=F_{m-1}+F_{m-2}+F_{m-3}.
\end{equation}
By Proposition~\ref{prop:cyclic-count},
\begin{equation}\label{eq:C123}
  C_n=F_{n-1}+2F_{n-2}+3F_{n-3}.
\end{equation}
The initial values are
\[
  C_0=3,\qquad C_1=1,\qquad C_2=3,
\]
and therefore $C_n=L_n$ from \eqref{eq:Lucas-def}.  This proves
Theorem~\ref{thm:main}.

A compact refinement is the full winding polynomial.

\begin{corollary}[Sector polynomial]\label{cor:sector-poly-r4}
For $n\ge4$,
\begin{equation}\label{eq:sector-poly-r4}
  W_{n,4}(u)=1+L_nu+L_nu^2+u^3.
\end{equation}
\end{corollary}

\subsection{Recurrences}

Substituting $L_n=(\Phi_n-2)/2$ into the Tribonacci recurrence gives the affine
order-three recurrence
\begin{equation}\label{eq:true-affine-recurrence}
  \boxed{\Phi_n=\Phi_{n-1}+\Phi_{n-2}+\Phi_{n-3}-4}
  \qquad(n\ge7),
\end{equation}
with $\Phi_4=24,\Phi_5=44,\Phi_6=80$.  Subtracting two consecutive instances
gives the homogeneous order-four recurrence
\begin{equation}\label{eq:true-homogeneous-recurrence}
  \boxed{\Phi_n=2\Phi_{n-1}-\Phi_{n-4}}
  \qquad(n\ge8).
\end{equation}

The Tribonacci--Lucas numbers are odd: the first three are odd, and each later
term is a sum of three odd terms.  Hence the theorem also proves
\begin{equation}\label{eq:divisibility4}
  4\mid\Phi_n\qquad(n\ge4).
\end{equation}
Define the corrected quarter sequence
\begin{equation}\label{eq:R-def}
  R_n:=\frac{\Phi_n}{4}=\frac{L_n+1}{2}.
\end{equation}
It begins
\begin{equation}\label{eq:R-values}
  6,11,20,36,66,121,222,408,750,1379,2536,4664,\dots
\end{equation}
and satisfies
\begin{equation}\label{eq:R-affine}
  R_n=R_{n-1}+R_{n-2}+R_{n-3}-1
  \qquad(n\ge7).
\end{equation}

\subsection{Generating functions}

Let the tail variable start at $n=4$.  From
\eqref{eq:true-homogeneous-recurrence},
\begin{align}
  \sum_{m\ge0}\Phi_{m+4}z^m
  &=\frac{24-4z-8z^2-16z^3}{1-2z+z^4} \label{eq:Phi-tail-gf}\\
  &=\frac{4(6-z-2z^2-4z^3)}{(1-z)(1-z-z^2-z^3)}.\nonumber
\end{align}
Dividing by four gives
\begin{equation}\label{eq:R-tail-gf}
  \boxed{
  \sum_{m\ge0}R_{m+4}z^m
  =\frac{6-z-2z^2-4z^3}{1-2z+z^4}.}
\end{equation}

If one includes the exceptional small-$n$ values
$1,1,2,6$ appearing in A004306, then
\begin{equation}\label{eq:A004306-full-gf}
 \sum_{n\ge0}A004306(n)z^n
 =\frac{1-z+2z^3+13z^4-3z^5-6z^6-10z^7}{1-2z+z^4},
\end{equation}
which agrees with the current A004306 entry \cite{oeisA004306}.

\section{The Mendelsohn formula sequence and A000382}\label{sec:ghost}

This section separates a valid recurrence identity from an invalid
combinatorial identification.

\subsection{The formula-defined sequence}

Mendelsohn defines the third-order Fibonacci sequence in his indexing as
\[
  1,2,4,7,13,24,44,81,149,\dots
\]
With our
$F_0=1,F_1=1,F_2=2,\dots$, his displayed equation (4), divided by four, gives
for $n\ge5$
\begin{equation}\label{eq:M-def}
  M_n=F_{n-1}+F_{n-3}+F_{n-4}+1,
\end{equation}
with the separately correct value $M_4=6$.  The sequence $M_n$ is exactly the
current A000382 tail and $4M_n$ is the current A000496 tail.

There is a harmless typographical issue in the preceding paragraph's display
in some text extractions of the 1961 scan; the sequence meant is the one
printed in the paper and used in \eqref{eq:M-def}.  The numerical evaluation
is unambiguous: $M_8=65,M_9=119,M_{10}=218$.

\begin{theorem}[Resolution of the A000382 recurrence]\label{thm:ghost-recurrence}
The formula-defined sequence $M_n$ satisfies
\begin{equation}\label{eq:ghost-affine}
  \boxed{M_n=M_{n-1}+M_{n-2}+M_{n-3}-2}
  \qquad(n\ge8).
\end{equation}
Its shifted generating function is
\begin{equation}\label{eq:ghost-gf}
  \boxed{
  \sum_{m\ge0}M_{m+4}z^m
  =\frac{6-z-2z^2-4z^3-z^4}{1-2z+z^4}.}
\end{equation}
Thus the recurrence marked conjectural in A000382 is a theorem for the
sequence defined by \eqref{eq:M-def}.  It is not the recurrence of the actual
quarter-count $R_n$.
\end{theorem}

\begin{proof}
The sequence $M_n-1$ is a fixed linear combination of three shifts of $F_n$,
so it obeys the homogeneous Tribonacci recurrence whenever all shifted
instances of \eqref{eq:M-def} are valid.  Therefore
\[
 M_n-1=(M_{n-1}-1)+(M_{n-2}-1)+(M_{n-3}-1),
\]
which is \eqref{eq:ghost-affine}.  The initial tail
$6,11,20,36,65$ then gives the numerator in \eqref{eq:ghost-gf} after
multiplication by $1-2z+z^4$.
\end{proof}

The homogeneous recurrence $M_n=2M_{n-1}-M_{n-4}$ begins only at $n=9$;
it fails at $n=8$ because the numerator in \eqref{eq:ghost-gf} contains a
nonzero $z^4$ term.  This one-step distinction is another useful check on
indexing.

\subsection{The exact correction kernel}

Subtracting \eqref{eq:ghost-gf} from \eqref{eq:R-tail-gf} gives the central
provenance identity.

\begin{theorem}[Exact correction to A000382 and A000496]\label{thm:correction}
For the current A000382 tail $M_n$ and the true quarter-count $R_n$,
\begin{equation}\label{eq:correction-gf}
  \boxed{
  \sum_{n\ge4}(R_n-M_n)z^{n-4}
  =\frac{z^4}{1-2z+z^4}
  =\frac{z^4}{(1-z)(1-z-z^2-z^3)}.}
\end{equation}
Consequently,
\begin{equation}\label{eq:correction-sum}
  R_n-M_n=
  \begin{cases}
    0,&4\le n\le7,\\[2mm]
    \displaystyle\sum_{j=0}^{n-8}F_j,&n\ge8.
  \end{cases}
\end{equation}
For A000496, the correction is four times this amount.
\end{theorem}

The first nonzero corrections to A000382 are therefore
\begin{equation}\label{eq:correction-values}
  1,2,4,8,15,28,52,96,177,\dots\qquad(n=8,9,10,\dots).
\end{equation}
They are not a finite transcription error: they grow exponentially at the same
Tribonacci rate as the two compared sequences.

A useful closed form for the cumulative composition numbers is
\begin{equation}\label{eq:cumulative-F}
  \sum_{j=0}^{m}F_j=\frac{F_{m+2}+F_m-1}{2}.
\end{equation}
It follows either by induction or by multiplying generating functions.

\subsection{What the 1961 source actually contains}

Mendelsohn's paper gives the correct definition of $\phi(n,r)$ and presents
its evaluation via permanents \cite{mendelsohn1961}.  For $r=4$, however,
its displayed equation (4) produces $260,476,872$ at $n=8,9,10$, while its
printed table gives $256,472,872$.  The permanent gives $264,484,888$.
The agreement through $n=7$ explains how the formula-defined sequence could
remain attached to the intended combinatorial description.

Plouffe's 1992/2009 collection explicitly describes its formulae as
computer-generated approximations and conjectures from the sequence database
available at the time \cite{plouffe2009}.  Its rational generating function
correctly captures the inherited A000382 terms; it cannot repair a prior error
in those terms without returning to the underlying definition.

\section{Exact roots, asymptotics, and nearest integers}\label{sec:asymptotics}

Let $\alpha$ be the real root of
\begin{equation}\label{eq:trib-polynomial}
  x^3-x^2-x-1=0,
\end{equation}
and let $\beta,\gamma$ be the two nonreal roots.  Numerically,
\begin{align}
 \alpha&=1.8392867552141611325518525646532866004\dots,\label{eq:alpha-num}\\
 \beta&=\rho\e^{\ii\theta},\qquad
 \gamma=\rho\e^{-\ii\theta},\nonumber\\
 \rho&=\alpha^{-1/2}
      =0.7373527057603276752017596505081211233\dots,\label{eq:rho-num}\\
 \theta&=2.1762335454918703984542598307236195688\dots.\label{eq:theta-num}
\end{align}
Indeed $\alpha\beta\gamma=1$, so $\abs\beta=\abs\gamma=\alpha^{-1/2}$, and
\begin{equation}\label{eq:theta-exact}
  \cos\theta=\frac{(1-\alpha)\sqrt\alpha}{2}.
\end{equation}

\begin{theorem}[Exact Binet form]\label{thm:Binet}
For every $n\ge4$,
\begin{align}
  \Phi_n
  &=2\bigl(1+\alpha^n+\beta^n+\gamma^n\bigr)\label{eq:Binet-complex}\\
  &=2\alpha^n+2+4\alpha^{-n/2}\cos(n\theta).\label{eq:Binet-real}
\end{align}
In particular,
\begin{equation}\label{eq:sharp-error}
  \abs{\Phi_n-(2\alpha^n+2)}\le4\alpha^{-n/2},
  \qquad
  \Phi_n\sim2\alpha^n.
\end{equation}
\end{theorem}

\begin{proof}
The power sum $\alpha^n+\beta^n+\gamma^n$ has initial values $3,1,3$ by
Newton's identities and obeys the Tribonacci recurrence.  Hence it is $L_n$.
Now apply Theorem~\ref{thm:main}.
\end{proof}

The result is stronger than a conventional asymptotic expansion: it is an
exact three-scale exponential representation.  The subdominant pair produces
a decaying oscillation rather than a monotone error.

\begin{corollary}[Nearest-integer formula]\label{cor:nearest}
For every $n\ge4$,
\begin{equation}\label{eq:nearest-quarter}
  \boxed{
  \frac{\Phi_n}{4}
  =\operatorname{nint}\!\left(\frac{\alpha^n+1}{2}\right),}
\end{equation}
where $\operatorname{nint}$ denotes the unique nearest integer.
\end{corollary}

\begin{proof}
By \eqref{eq:Binet-real}, the difference between the two sides before rounding
is $\rho^n\cos(n\theta)$, whose absolute value is at most $\rho^n<1/2$ for
$n\ge3$.
\end{proof}

\subsection{A convergent logarithmic transseries}

Factoring the dominant term in \eqref{eq:Binet-real} gives
\begin{equation}\label{eq:factor-Phi}
 \Phi_n=2\alpha^n
 \left(1+\alpha^{-n}+2\alpha^{-3n/2}\cos(n\theta)\right).
\end{equation}
Therefore, for all sufficiently large $n$ (and in fact already in the stated
range after direct checking),
\begin{equation}\label{eq:log-all-orders}
 \log\Phi_n
 =n\log\alpha+\log2+
 \sum_{m\ge1}\frac{(-1)^{m+1}}m
 \left(\alpha^{-n}+2\alpha^{-3n/2}\cos(n\theta)\right)^m.
\end{equation}
This is a convergent all-orders expansion, not merely a formal one.  Its first
layers are
\begin{align}
 \log\Phi_n={}&n\log\alpha+\log2
 +\alpha^{-n}+2\alpha^{-3n/2}\cos(n\theta)
 -\frac12\alpha^{-2n}\nonumber\\
 &-2\alpha^{-5n/2}\cos(n\theta)
 +\left(\frac13-2\cos^2(n\theta)\right)\alpha^{-3n}
 +\bigO(\alpha^{-7n/2}).\label{eq:log-expanded}
\end{align}
The half-integer exponential scales arise from the modulus
$\abs\beta=\alpha^{-1/2}$.

\subsection{Asymptotics of the erroneous formula sequence}

The composition numbers have the exact root expansion
\begin{equation}\label{eq:F-root-expansion}
  F_n=\sum_{\lambda\in\{\alpha,\beta,\gamma\}}
  \frac{\lambda^{n+3}}{\lambda^2+2\lambda+3}.
\end{equation}
Substitution into \eqref{eq:M-def} yields, for $n\ge5$,
\begin{equation}\label{eq:M-root-expansion}
 M_n=1+\sum_{\lambda\in\{\alpha,\beta,\gamma\}}
 \frac{\lambda^{n-1}(\lambda^2+2\lambda+2)}
      {\lambda^2+2\lambda+3}.
\end{equation}
Thus
\begin{equation}\label{eq:M-asymptotic}
 M_n\sim
 \frac{\alpha^2+2\alpha+2}
      {\alpha(\alpha^2+2\alpha+3)}\,\alpha^n
 =0.4896527010215867240\dots\,\alpha^n.
\end{equation}
Since $R_n\sim\frac12\alpha^n$, the limiting relative undercount is
\begin{equation}\label{eq:relative-undercount}
  1-\lim_{n\to\infty}\frac{M_n}{R_n}
  =0.0206945979568265519\dots.
\end{equation}
So the discrepancy is small enough to be visually deceptive, but it does not
vanish relatively.

\section{Inverting the exact asymptotic transseries}\label{sec:inverse}

The exact root formula permits a precise continuous inverse and, through it,
a discrete threshold formula.

Define the real interpolation
\begin{equation}\label{eq:continuous-Phi}
  \widetilde\Phi(x)
  :=2+2\alpha^x+4\alpha^{-x/2}\cos(\theta x),
  \qquad x\ge4.
\end{equation}
It agrees with $\Phi_n$ at every integer $n\ge4$.

\begin{lemma}[Monotonicity of the interpolation]\label{lem:interpolation-monotone}
The function $\widetilde\Phi$ is strictly increasing on $[4,\infty)$.
\end{lemma}

\begin{proof}
Put $L=\log\alpha$ and
$B=\sqrt{(L/2)^2+\theta^2}$.  Then
\[
 \widetilde\Phi'(x)
 \ge 2L\alpha^x-4B\alpha^{-x/2}.
\]
The polynomial in \eqref{eq:trib-polynomial} is increasing for $x>1$ and
changes sign between $1.8$ and $2$, so $1.8<\alpha<2$.  Hence
$L>0.58$, while $\theta<\pi$ and
$B<\sqrt{(\log 2/2)^2+\pi^2}<3.17$.  At $x=4$ the displayed lower bound is
therefore greater than
\[
  2(0.58)(1.8)^4-4(3.17)(1.8)^{-2}>8.2.
\]
The ratio of the possible negative term to the positive term decreases by a
factor $\alpha^{-3/2}$ per unit increase of $x$, so the derivative remains
positive.
\end{proof}

For $y\ge24$, let $x(y)$ be the unique real solution of
$\widetilde\Phi(x)=y$.  Set
\begin{align}\label{eq:inverse-variables}
  Y&:=\frac{y-2}{2},\qquad L:=\log\alpha,\\
  \omega&:=\frac{\theta}{L}
  =3.5712382678640038805657341982470\dots,
  \qquad \psi:=\omega\log Y.\nonumber
\end{align}

\begin{theorem}[Inverse log-periodic transseries]\label{thm:inverse}
As $y\to\infty$,
\begin{align}
 x(y)={}&\frac{\log Y}{L}
 -\frac{2\cos\psi}{L}Y^{-3/2}\nonumber\\
 &-\left(
   \frac{4\cos^2\psi}{L}
   +\frac{4\theta\cos\psi\sin\psi}{L^2}
   \right)Y^{-3}
 +\bigO(Y^{-9/2}).\label{eq:inverse-expansion}
\end{align}
More generally, there is an all-orders expansion
\begin{equation}\label{eq:inverse-all-orders}
 x(y)=\frac{\log Y}{L}
      +\sum_{m\ge1}a_m(\psi)Y^{-3m/2},
\end{equation}
where each $a_m$ is a real-analytic $2\pi$-periodic function computable
recursively.
\end{theorem}

\begin{proof}
Write
\[
  x=\frac{\log Y}{L}+\delta,
  \qquad t=Y^{-3/2}.
\]
Dividing the equation $Y=\alpha^x+2\alpha^{-x/2}\cos(\theta x)$ by $Y$ gives
\begin{equation}\label{eq:inverse-implicit}
  1=\e^{L\delta}
    +2t\e^{-L\delta/2}\cos(\psi+\theta\delta).
\end{equation}
At $t=0$ the solution is $\delta=0$, and the derivative of the right side
with respect to $\delta$ is $L\ne0$.  The analytic implicit-function theorem
therefore produces a convergent power series
$\delta=\sum_{m\ge1}a_m(\psi)t^m$ for small $t$, uniformly in the phase.

Substitute $\delta=a_1t+a_2t^2+\bigO(t^3)$ into
\eqref{eq:inverse-implicit}.  Comparing the coefficient of $t$ gives
\[
  La_1+2\cos\psi=0.
\]
Comparing the coefficient of $t^2$ then gives
\[
  La_2+4\cos^2\psi
  +\frac{4\theta}{L}\cos\psi\sin\psi=0.
\]
These are exactly the two displayed correction terms.
\end{proof}

\begin{corollary}[Discrete inverse]\label{cor:discrete-inverse}
For $y\ge24$, define
\[
 N(y):=\min\{n\ge4:\Phi_n\ge y\}.
\]
Then
\begin{equation}\label{eq:discrete-inverse}
  \boxed{N(y)=\ceil{x(y)}.}
\end{equation}
Consequently, inserting \eqref{eq:inverse-expansion} gives an asymptotically
sharp threshold formula with a first correction of order $y^{-3/2}$ and a
log-periodic phase.
\end{corollary}

\begin{proof}
Lemma~\ref{lem:interpolation-monotone} shows that
$\widetilde\Phi(n)\ge y$ if and only if $n\ge x(y)$.  Since
$\widetilde\Phi(n)=\Phi_n$, the least such integer is $\ceil{x(y)}$.
\end{proof}

\section{Independent exact verification}\label{sec:verification}

The supplied file \texttt{verify.py} performs an independent calculation using
only the Python standard library.  It does not use the recurrence proved in
this article to compute the permanent.

For row $i$, a dynamic-programming state records:
\begin{enumerate}[label=(\alph*),leftmargin=2.3em]
\item the bit mask of already used target columns;
\item the exact integer sum $d_0+\cdots+d_i$.
\end{enumerate}
Each transition chooses one of the four allowed targets
$i,i+1,i+2,i+3\pmod n$, provided that target has not been used.  At the final
row, the displacement sum is divided by $n$ to recover the winding sector.
All arithmetic is integral.

The computation through $n=16$ returns:
\begin{center}
\begin{tabular}{r@{\quad}r@{\quad}r@{\quad}r@{\quad}r@{\quad}r}
\toprule
$n$ & $q=0$ & $q=1$ & $q=2$ & $q=3$ & total\\
\midrule
4  &1&11&11&1&24\\
5  &1&21&21&1&44\\
6  &1&39&39&1&80\\
7  &1&71&71&1&144\\
8  &1&131&131&1&264\\
9  &1&241&241&1&484\\
10 &1&443&443&1&888\\
11 &1&815&815&1&1632\\
12 &1&1499&1499&1&3000\\
13 &1&2757&2757&1&5516\\
14 &1&5071&5071&1&10144\\
15 &1&9327&9327&1&18656\\
16 &1&17155&17155&1&34312\\
\bottomrule
\end{tabular}
\end{center}

The script also checks:
\begin{itemize}[leftmargin=2em]
\item both exact recurrences for the true count;
\item the A000382 recurrence for the formula-defined sequence;
\item both shifted generating-function numerators;
\item the correction kernel coefficient by coefficient;
\item the nearest-integer identity through $n=100$ using an 80-digit Newton
      solve for $\alpha$.
\end{itemize}
The captured output is included as \texttt{verification\_output.txt}.

\section{Proposed OEIS corrections}\label{sec:oeis-corrections}

The following are mathematical recommendations, not submissions.  Attribution,
formatting, and historical priority should be reviewed by OEIS editors.

\subsection{A000382}

\begin{enumerate}[leftmargin=2.2em]
\item Remove the comment ``The fourth column of A008305, divided by 4.''  It is
      false beginning at $n=8$.
\item Either rename the entry as the sequence defined by Mendelsohn's displayed
      formula (4), divided by $4$, or mark the historical provenance explicitly.
\item Replace ``conjectured'' on the affine recurrence by a proof reference to
      Theorem~\ref{thm:ghost-recurrence}, provided the formula-defined meaning
      is adopted.
\item Label the rational function as the shifted tail generating function
      $\sum_{m\ge0}a(m+4)z^m$.
\item Cross-reference the corrected quarter-count $R_n=(A001644(n)+1)/2$ for
      $n\ge4$ and A004306 for the undivided count.
\end{enumerate}

\subsection{A000496}

\begin{enumerate}[leftmargin=2.2em]
\item Add that the terms from $n=8$ onward are not the permanent
      $\per(I+P+P^2+P^3)$ and are not the fourth column of A008305.
\item Clarify that the displayed rational function is a tail generating
      function, not the ordinary generating function beginning at offset $0$.
\item Retain the historical sequence if desired, but distinguish it from
      A004306.
\end{enumerate}

\subsection{A008305}

In the cross-reference line for the fourth column, remove ``$4*$A000382'' and
retain A004306.  A direct comment may state that for $n\ge4$
\begin{equation}\label{eq:oeis-direct-comment}
  T(n,4)=2(A001644(n)+1),
\end{equation}
with sector sizes $(1,A001644(n),A001644(n),1)$.

\subsection{A004306}

The current entry already contains the correct recurrence and the equivalent
Tribonacci formula.  Useful additions would be:
\begin{align*}
 A004306(n)&=2\left(1+\alpha^n+\beta^n+\gamma^n\right),\qquad n\ge4,\\
 \frac{A004306(n)}4
 &=\operatorname{nint}\left(\frac{\alpha^n+1}{2}\right),\qquad n\ge4,
\end{align*}
and the winding-sector tiling interpretation from this report.

A plain-text draft suitable for editing is included as
\texttt{oeis\_corrections.txt}.

\section{Further research questions and topics}\label{sec:further}

The winding decomposition is complete for $r=4$ but exposes a larger program.
The following problems are intended to be concrete enough for computation,
proof, or formalization.

\begin{question}[Inner winding sectors for $r\ge5$]
For fixed $r$, determine the numbers
\[
  N_{n,r}(q):=\#\{\sigma:q(\sigma)=q\}
\]
for $2\le q\le r-3$.  The constant-cut law turns this into the enumeration of
$q$-fold covers of a labeled cycle by intervals of lengths at most $r-1$ with
the additional bijectivity constraints at endpoints.  Find a canonical local
object---multitilings, cylindric paths, or a finite-state model---that makes
those constraints transparent.
\end{question}

\begin{question}[Sector polynomial shape]
The involution proves that $W_{n,r}(u)$ is palindromic.  Is it unimodal or
log-concave for all $n\ge r$?  Is there a stability, real-rootedness, or
$\gamma$-positivity phenomenon after a natural normalization?  Small cases
should be classified before making a universal conjecture.
\end{question}

\begin{question}[Minimal recurrences at fixed width]
Classical transfer-matrix theory implies rational generating functions in $n$
for each fixed $r$, but the minimal denominator and its factorization can be
subtle.  Determine the minimal recurrence for each individual sector and for
the total $\Phi_{n,r}$, and explain cancellations bijectively.
\end{question}

\begin{question}[Dominant winding sector]
For $r\ge5$, which sector controls the exponential growth of $\Phi_{n,r}$?
Does the maximizing sector concentrate near $(r-1)/2$?  Obtain the exponential
rate and leading constant from a Perron--Frobenius operator, then compare with
Mendelsohn's general asymptotic claim.
\end{question}

\begin{question}[Weighted displacement enumerators]
Study
\[
 \per\!\left(w_0I+w_1P+\cdots+w_{r-1}P^{r-1}\right)
\]
sector by sector.  For $r=4$, the winding-one sector becomes a weighted cyclic
composition polynomial.  Determine whether the full four-sector formula has a
useful multivariate factorization and whether it yields refined statistics on
fixed points and jump lengths.
\end{question}

\begin{question}[Arbitrary displacement sets]
Replace $\{0,1,\dots,r-1\}$ by an arbitrary subset $S\subset\Zn$.  Characterize
those $S$ for which an extreme nonzero winding sector is a tiling language, and
those for which the permanent admits a low-order recurrence independent of
$n$.
\end{question}

\begin{question}[Other finite groups]
Mendelsohn formulated displacements in abelian groups but specialized much of
the computation to cyclic groups.  Develop an analogue of the cut-flow
argument on products of cycles, where cuts become codimension-one tori and the
winding number becomes a homology class.  Even the first nontrivial case
$\Z_m\times\Z_n$ may reveal a useful higher-dimensional tiling model.
\end{question}

\begin{question}[A natural model for the historical ``ghost'' sequence]
The A000382 tail has a perfectly valid positive recurrence and differs from the
true quarter-count by cumulative Tribonacci numbers.  Is there a natural class
of cyclic tilings with a marked defect, boundary condition, or forbidden wrap
that is counted by $M_n$?  A satisfactory model would preserve the historical
sequence while preventing future confusion about its meaning.
\end{question}

\begin{question}[Divisibility without recurrence]
The formula proves $4\mid\Phi_n$.  Find a direct free action or sign-reversing
argument that partitions the allowed permutations into quadruples.  Such an
action might generalize to divisibility properties of $\Phi_{n,r}$.
\end{question}

\begin{question}[Inverse transseries for general rational recurrences]
The inverse in Section~\ref{sec:inverse} has periodic coefficients because the
subdominant roots are a complex conjugate pair.  Develop an automated theorem
that takes a positive linear recurrence with a simple dominant root and
returns the inverse transseries, including resonances among subdominant root
moduli and phases.
\end{question}

\begin{question}[Formal verification]
Formalize the constant-cut lemma, tiling bijection, and recurrence in Lean.
The proof requires only finite sums, permutations of \texttt{Fin n}, modular
arithmetic, and cyclic interval bookkeeping, making it a realistic addition to
the ProveIt formalization program.
\end{question}

\begin{question}[Automated provenance audits]
Build a reproducible tool that compares an OEIS definition against its b-file,
formulae, cited tables, and independently computed terms.  The present example
shows why matching a long recurrence to inherited terms is not a substitute
for recomputing the defining combinatorial object.
\end{question}

\section{Conclusion}

The four-displacement problem is governed by a simple topological statistic.
The sum of clockwise displacements, divided by $n$, is an integer winding
number, and local flow conservation forces every cut to have the same
multiplicity.  At winding one this turns the permutation into a cyclic tiling;
for four allowed displacements, symmetry accounts for every remaining sector.
The resulting exact count is
\[
  \Phi_n=2(L_n+1),
\]
not the sequence currently attached to A000382 and A000496.

The historical error is unusually instructive.  A published formula, a table
in the same paper, and a later machine-guessed generating function all agree
long enough to look persuasive, yet none of them replaces evaluation of the
original permanent.  Once the combinatorial structure is recovered, the
correct recurrence, generating function, Binet formula, asymptotics, and
inverse expansion all follow with little ambiguity.

\appendix

\section{Proof details for the generating functions}

For reference, multiplying a tail sequence $a_0,a_1,\dots$ by
$1-2z+z^4$ produces coefficients
\[
 a_m-2a_{m-1}+a_{m-4},
\]
where missing terms are zero.  For the true quarter-tail
$6,11,20,36,66,\dots$, these coefficients are
\[
 6,-1,-2,-4,0,0,\dots,
\]
which proves \eqref{eq:R-tail-gf}.  For the A000382 tail
$6,11,20,36,65,\dots$, they are
\[
 6,-1,-2,-4,-1,0,0,\dots,
\]
which proves \eqref{eq:ghost-gf}.  Their difference is exactly $z^4$ over the
common denominator.

\section{Reproduction instructions}

The archive contains:
\begin{description}[leftmargin=2.8cm,style=nextline]
\item[\texttt{article.tex}] this source file;
\item[\texttt{article.pdf}] the compiled report;
\item[\texttt{verify.py}] independent exact enumeration and checks;
\item[\texttt{verification\_output.txt}] captured verifier output;
\item[\texttt{oeis\_corrections.txt}] editable proposed OEIS notes;
\item[\texttt{README.md}] a concise file guide and build command;
\item[\texttt{Makefile}] convenience targets for verification and compilation.
\end{description}

Run
\begin{lstlisting}[language=bash]
python3 verify.py
pdflatex -interaction=nonstopmode article.tex
pdflatex -interaction=nonstopmode article.tex
\end{lstlisting}
or use the PDF build command documented in the README.

\begin{thebibliography}{99}

\bibitem{mendelsohn1961}
N.~S. Mendelsohn,
\emph{Permutations with confined displacements},
Canadian Mathematical Bulletin \textbf{4} (1961), 29--38.
\url{https://doi.org/10.4153/CMB-1961-005-4}.

\bibitem{metropolis1969}
N. Metropolis, M.~L. Stein, and P.~R. Stein,
\emph{Permanents of cyclic $(0,1)$ matrices},
Journal of Combinatorial Theory \textbf{7} (1969), 291--321.
\url{https://doi.org/10.1016/S0021-9800(69)80058-X}.

\bibitem{beker1987}
H. Beker and C. Mitchell,
\emph{Permutations with restricted displacement},
SIAM Journal on Algebraic Discrete Methods \textbf{8} (1987), 338--363.
\url{https://doi.org/10.1137/0608029}.

\bibitem{plouffe2009}
S. Plouffe,
\emph{Approximations of generating functions and a few conjectures},
arXiv:0911.4975, 2009; based on a 1992 dissertation and formula collection.
\url{https://arxiv.org/abs/0911.4975}.

\bibitem{oeisA000382}
OEIS Foundation Inc.,
\emph{A000382: Restricted permutations}, snapshot checked October 1, 2026.
\url{https://oeis.org/A000382}.

\bibitem{oeisA000496}
OEIS Foundation Inc.,
\emph{A000496: Restricted permutations}, snapshot checked October 1, 2026.
\url{https://oeis.org/A000496}.

\bibitem{oeisA004306}
OEIS Foundation Inc.,
\emph{A004306: Rook polynomials; fourth column of A008305},
snapshot checked October 1, 2026.
\url{https://oeis.org/A004306}.

\bibitem{oeisA008305}
OEIS Foundation Inc.,
\emph{A008305: permutations with cyclic restricted displacements},
snapshot checked October 1, 2026.
\url{https://oeis.org/A008305}.

\bibitem{oeisA001644}
OEIS Foundation Inc.,
\emph{A001644: Tribonacci--Lucas numbers},
snapshot checked October 1, 2026.
\url{https://oeis.org/A001644}.

\bibitem{proveit}
V. Reshetnikov,
\emph{ProveIt}, GitHub repository.
\url{https://github.com/VladimirReshetnikov/ProveIt}.

\end{thebibliography}

\end{document}
